\section{Sexta iteración de la Fase 5: Creación del GWC}\label{sec:F5I6}

Con el módulo preparado, llega el nuevo código: generar el
\textbf{GWC}, el cartucho de ejecución definido en el formato de la
segunda iteración definida en la sección \ref{sec:F5I2}. Todo se
incluye en \texttt{BaseCartucho.js} que ya se preparó para el
traspaso, con dos nuevas funciones: la \textbf{lectura} y la
\textbf{grabación}.

\subsection{El botón «Generar GWI+GWC»}
\label{sec:F5I6Boton}

La generación se lanza desde el editor gráfico, que tiene toda la
información, con un botón nuevo junto al de «Verificar y Guardar».
El botón del GWI solo SE MANTIENE: poder dejar un cartucho a medias
y seguir luego es importante. El nuevo botón encadena los procesos
sobre las mismas listas en memoria, con la verificación completa del
diseñador al frente, bloqueante:

\begin{enumerate}	
	\item \textbf{Generación del GWC}: ensamblado de los bloques,
	cabecera con su suma de verificación y descarga
	(\texttt{cartucho\_gwc.gwc}).	
	
	\item \textbf{Guardado del GWI} con el campo \texttt{codigo}
	igual a la \textbf{suma de verificación del GWC} recién
	generado, y su descarga (\texttt{cartucho\_gwi.zip}).
\end{enumerate}

La razón del encadenamiento es la atomicidad: si la validación
falla, no se genera ninguno de los dos; si tiene éxito, ambos salen
de la misma información de la memoria y son \textbf{siempre copia
	uno del otro}. No existe la situación de un GWC huérfano de un GWI
más nuevo. Los avisos de la barra de estado acompañan cada paso
(«Generando GWC…», «Guardando GWI…»), y cualquier fallo es visible
con su mensaje, como en todo el proyecto.

La verificación de los topes acompaña este reparto, y lo hace como
función independiente: \texttt{validarTopes(listas)} —en el
\texttt{BaseCartucho.js}— devuelve \textbf{cierto o falso}, y es el
punto de llamada quien decide su tratamiento: tras descomponer los
segmentos y al guardar el GWI su fallo es solo un \textbf{aviso}
—el diseñador sabe que se ha pasado del tope y decide si seguir—, y
al pulsar «Generar GWI+GWC» es \textbf{bloqueante}: no da el error
cuando uno piensa que está todo bien, sino en el único momento en
que el GWC ya no puede salir correcto.

Lo mismo ocurre con la verificación del diseñador: la coherencia de
enlaces, contadores y tipos que se lanzaba antes de guardar el GWI pasa a
su \textbf{propia función independiente}, que los dos botones de
guardar llaman \textbf{siempre antes de empezar} su proceso —un
solo código de verificación, sin copias— y que es bloqueante en
ambos: si la verificación falla, ni el GWI ni el par se guardan.

\subsection{La grabación: \texttt{generarGWC(listas)}}\label{sec:F5I6Grabacion}

Para poder leer los datos, primero hay que generarlos, razón por la
que primero abordo el tema de la grabación y luego el de la lectura.
Se definen estas funciones:

\begin{table}[H]
	\centering
	\begin{tabularx}{\textwidth}{@{}l X@{}}
		\toprule
		\textbf{Función} & \textbf{Descripción} \\
		\midrule
		
		\texttt{leerGWC(bytes)} & La lectura: valida la
		cabecera y la suma de verificación, recorre los bloques y
		reconstruye las listas en memoria con las mismas formas que
		produce el editor \\
		
		\texttt{generarGWC(listas)} & La grabación: valida los
		topes, escala, ensambla los bloques en orden, escribe la
		cabecera con la suma de verificación y devuelve los bytes
		junto a la estadística de tamaños \\
		
		\texttt{validarTopes(listas)} & Comprueba que sprites,
		segmentos y secuencias no superan los 255 elementos
		direccionables con índices de 1 byte: en el editor es un
		\textbf{aviso} —tras descomponer los segmentos y al guardar
		el GWI—, y en el botón «Generar GWI+GWC» es
		\textbf{bloqueante}: no se genera nada \\
		
		\texttt{escalar(sprite, pantalla)} & Convierte un sprite en
		su máscara de 1 bit escalada a la pantalla declarada, con el
		muestreo por vecino más cercano y escala por eje del
		simulador: lo que viaja es lo que el simulador pinta al vuelo \\
		
		\texttt{indiceDe(id)} & Traduce los \texttt{seg\_NNN} del
		GWI a índices del GWC, con el \textbf{nulo} para los enlaces
		ausentes; los campos \emph{salvadoContador} y
		\emph{muerteContador} viajan con el contador correcto o
		nulos \\
		
		\texttt{bloque(tipo, contenido)} & Escribe la estructura
		común de todos los bloques: tipo, tamaño del contenido de 4
		bytes y contenido \\
		
		\texttt{estadistica(bytes)} & Presenta el tamaño final del
		GWC generado, con la descomposición por bloques \\
		\bottomrule
	\end{tabularx}
	\caption{Las funciones del GWC en \texttt{BaseCartucho.js}: leer antes que grabar, como los otros dos formatos.}
	\label{tab:funcionesGWC}
\end{table}

El generador no usa el azar en ninguna parte: recorre las listas
siempre en el mismo orden y escribe siempre lo mismo. La generación
es \textbf{determinista}: el mismo GWI produce el mismo GWC byte a
byte, con la misma suma de verificación —el equivalente del
compilador reproducible, y la base de su prueba de regresión.

El descarte también vive aquí: los identificadores legibles, los
fantasmas tipo 99, los enlaces inversos calculados y toda la
sección de autoría no viajan; el generador los descarta al vuelo,
un segmento que nadie enciende ni pinta no existe en el GWC.

\subsection{La lectura: \texttt{leerGWC(bytes)}}\label{sec:F5I6Lectura}

La función de lectura es pura de bytes a listas, sin DOM, sin ZIP y
sin imágenes PNG: recorre la cabecera —identificador «GWC»,
versión, tamaño del fichero y suma de verificación—, la
\textbf{valida}, y recorre los bloques en su orden, reconstruyendo
las listas en memoria con las mismas formas que produce el editor:
los índices binarios vuelven a ser referencias de segmento, las
máscaras de 1 bit vuelven a ser imágenes en memoria, y los campos
de configuración vuelven al puente como si vinieran de un GWI.

El bloque que no existe no se echa en falta: el registro de
secuencias ausente deja la Secuencia a \texttt{null} y el bloque de
fondo ausente deja el color de fondo declarado en el Bloque 0.
Cualquier desviación —identificador equivocado, versión no
soportada, tamaño que no cuadra con el fichero, suma de
verificación que no coincide— es un error explícito con su
mensaje, como en todo el proyecto: un cartucho corrupto no se
carga a medias, y el cartucho anterior se conserva.

Esta función queda en la \textbf{sección delimitada y
	autocontenida} del módulo: sin dependencias del GWB, del ZIP ni del
editor, para servir —con los cambios oportunos— como lector del
firmware del ESP32 de la siguiente fase.

\subsection{El código del par: la suma de verificación}\label{sec:F5I6Codigo}

Para asegurarse de que un GWI y un GWC salieron juntos de la misma
memoria, el GWI lleva un campo \texttt{codigo} que, al generar el
par, vale exactamente \textbf{la suma de verificación del GWC}
—el mismo dato de 16 bits que el GWC lleva en su cabecera. El
formato del GWC no gana ningún byte: su identificador de par ya
existe, es su propia firma.

El reparto es el mismo de siempre:

\begin{itemize}	
	\item \textbf{«Verificar y Guardar»} deja el GWI
	\textbf{en blanco}: guardado a solas, un GWI no está emparejado
	con nada.	
	\item \textbf{«Generar GWI+GWC»} escribe primero el GWC y
	traslada su suma de verificación al campo \texttt{codigo} del
	GWI: ambos ficheros quedan sellados con el mismo dato.
\end{itemize}

Comprobar que van de la mano es leer dos valores: el
\texttt{codigo} del GWI y la suma de verificación de la cabecera
del GWC. Que un par coincida por casualidad exigiría que dos
cartuchos distintos compartieran suma de verificación —la
probabilidad de que los cambios se anulen es despreciable en
nuestro caso. Y la determinación del código no rompe nada: como la
generación es determinista, un mismo GWI produce siempre el mismo
GWC y por tanto el mismo código —la ida y vuelta sigue siendo
válida byte a byte.

\subsection{La carga en el simulador y la verificación}\label{sec:F5I6Carga}

El generador gráfico permite cargar un cartucho GWC con
\texttt{leerGWC}: se verifica que su contenido es válido, que
visualmente carga lo mismo que un GWI, que los segmentos conservan
sus propiedades y que el simulador ejecuta el mismo juego.

Y esto permite la verificación más exigente, el \textbf{ciclo de
	ida y vuelta}:

\begin{enumerate}	
	\item Se carga un GWI.	
	\item Sin tocar nada se genera el GWI+GWC.	
	\item Se carga el GWC generado.	
	\item Sin tocar nada se genera de nuevo el GWI+GWC.
\end{enumerate}

Los dos GWC —el del punto 2 y el del punto 4— deben ser
\textbf{idénticos byte a byte}, muy sencillo de validar con los
comandos estándar del sistema. El GWI regenerado es
\textbf{equivalente en su contenido de juego}: la autoría y los
gráficos originales no viajan en el GWC, así que los GWI no pueden
ser idénticos, pero las listas que producen deben ser las mismas
—y el campo \texttt{codigo}, que es la firma del GWC, también debe
salir el mismo en ambos guardados.

\subsection{Simulación}\label{sec:F5I6Simulacion}

El generador y el lector son puro pasar listas en memoria a bytes
y al contrario, sin DOM, así que la batería de pruebas los ejercita
enteros con listas sintéticas y con el GWI real del Parachute:

\begin{itemize}	
	\item \textbf{Determinismo:} dos generaciones del mismo GWI
	producen bytes idénticos y la misma suma de verificación.	
	
	\item \textbf{Estructura:} los conteos de la cabecera cuadran
	con las tablas; toda dirección apunta dentro del fichero; la
	suma de verificación recomputada coincide; los punteros de los
	bloques opcionales ausentes son nulos.	
	
	\item \textbf{Contenido:} los enlaces del registro se resuelven
	contra los índices reales; ningún enlace válido apunta al nulo
	ni al revés; \emph{salvadoContador} y \emph{muerteContador}
	viajan correctos o nulos; no viajan fantasmas tipo 99 ni
	absorbidos.	
	
	\item \textbf{Escalado:} posiciones y máscaras coinciden a
	píxel con lo que el simulador pinta al vuelo con la misma
	pantalla declarada.	
	
	\item \textbf{Topes:} un GWI con más de 255 segmentos produce el
	aviso tras descomponer y al guardar el GWI, y el error
	bloqueante al pulsar «Generar GWI+GWC», sin generar ningún
	fichero.	
	
	\item \textbf{Código del par:} «Verificar y Guardar» produce un
	GWI sin código; «Generar GWI+GWC» pone en el campo
	\texttt{codigo} del GWI la misma suma de verificación que la
	cabecera del GWC; el código sobrevive a la ida y vuelta.	
	
	\item \textbf{Ida y vuelta:} cargar el GWC generado y volver a
	generar produce el mismo GWC byte a byte; el GWI regenerado es
	equivalente en su contenido de juego.	
	
	\item \textbf{Lectura corrupta:} un byte cambiado en la suma de
	verificación, un identificador equivocado o un tamaño que no
	cuadran dan su error explícito y el cartucho anterior se
	conserva.
\end{itemize}

Y en navegador real, con el cartucho del Parachute: el botón
produce ambos ficheros seguidos; un cartucho con un error de
verificación no produce ninguno; las cifras de la estadística se
presentan en la barra; la inspección del GWC (volcado hexadecimal)
muestra el identificador «GWC», su versión y el recorrido de
bloques en orden.

\subsection{La estadística de tamaños}\label{sec:F5I6Estadistica}

Al terminar, la generación presenta en la barra de estado y en
detalle bajo GW\_DEBUG una única cifra: el \textbf{tamaño final del
	GWC}, con la descomposición por bloques. El ahorro de la
deduplicación de sprites y la comparación frente al ZIP del GWI son
un tema aparte, más propio de la presentación final del proyecto, y
queda anotado como posible ampliación.

El Parachute real es el caso de estudio: su tamaño final en el GWC
—recogido en la memoria al cierre de la fase— ilustra la distancia
entre un formato de autoría (legible, editable, con historial e
imágenes PNG) y uno de ejecución (binario, mínimo, listo para el
microcontrolador). Los datos finales para la pantalla declarada de $320 \times 240$ han sido:


% ---------------------------------------------------------------
% Cifras del cartucho Parachute (F5-I7)
% ---------------------------------------------------------------
\begin{table}[H]
	\centering	
	\begin{tabular}{lr}
		\toprule
		Concepto & Valor \\
		\midrule
		Segmentos en el editor & 139 \\
		Imágenes en el editor & 125 \\
		GWI (JSON + PNG, aproximado) & 285\,558 bytes \\
		PNG de segmentos & 129\,352 bytes \\
		GWC generado & 7\,342 bytes \\
		\bottomrule
	\end{tabular}
	\caption{Cifras del cartucho Parachute en el editor y en la generación}
	\label{tab:F5I7CifrasParachute}
\end{table}

\begin{table}[H]
	\centering
	\begin{tabular}{llrr}
		\toprule
		Tipo & Bloque & Tamaño & Contenido \\
		\midrule
		- & Cabecera & 14 & \texttt{GWC1\#0BAD\#} + tamaño \\
		0 & Datos & 40 & estilo, pantalla, conteos, direcciones \\
		1 & Configuración & 21 & pausas, aceleración, rutas, fin \\
		2 & Sprites & 5\,523 & 52 máscaras \\
		3 & Segmentos & 1\,733 & 54 registros de 32 bytes \\
		4 & Secuencias & 11 & 1 secuencia de 5 segmentos \\
		5 & Fondo & -- & reservado (nulo) \\
		\midrule
		&Total & 7\,342 & \\
		\bottomrule
	\end{tabular}
	\caption{Desglose del GWC del Parachute por bloques}
	\label{tab:F5I7BloquesParachute}
\end{table}

\begin{table}[H]
	\centering
	\begin{tabular}{lr}
		\toprule
		Concepto & Valor \\
		\midrule
		Imágenes deduplicadas (compartidas en el editor) & 8 \\
		Ahorro en el GWI & 3\,829 bytes ($\approx 1'1\%$) \\
		Ahorro en el GWC & 52 bytes ($\approx 1'1\%$) \\
		\bottomrule
	\end{tabular}
	\caption{Ahorro por deduplicación de sprites en el Parachute}
	\label{tab:F5I7DedupParachute}
\end{table}

\subsection{Deduplicación aproximada}
\label{sec:F5I6Estadistica}


Si buscamos imágenes que sean casi iguales aunque no idénticas, analizando el porcentaje de diferencia entre ambas, para intentar fusionarlas y reducir un poco mas el tamaño, vemos que hay 3 pares «casi iguales» que la deduplicación exacta no fusiona:

\begin{itemize}
	\item par 24×20: 8 píxeles distintos (1,67\% de diferencia entre ambos)
	\item par 13×10: 3 píxeles distintos (2,31\% de diferencia entre ambos) 
	\item par 4×10: 3 píxeles distintos (7,5\% de diferencia entre ambos)
\end{itemize}

Con umbral del dos y medio por ciento entrarían los dos primeros, con un ahorro adicional de unos 88-102 bytes de GWC (~1,3\%). Es decir, usar una deduplicación de imágenes que sea aproximada añadiría poco sobre la actual que busca imágenes exactamente iguales, \underline{por lo que creo que no merece la pena implementarla}.
